\documentclass[12pt]{article}
\usepackage{amsmath, amssymb, amsthm}
\usepackage{geometry}
\usepackage{hyperref}
\usepackage{booktabs}
\usepackage{longtable}
\usepackage{graphicx}
\usepackage{titlesec}
\numberwithin{equation}{section}
\geometry{a4paper, margin=2.5cm}

% Stile elegante per i capitoli
\titleformat{\section}{\large\bfseries\centering}{}{0em}{}
\titlespacing*{\section}{0pt}{12pt}{6pt}
\titleformat{name=\section,numberless}{\large\bfseries\centering}{}{0em}{}
\titlespacing*{name=\section,numberless}{0pt}{12pt}{6pt}

\hypersetup{
    colorlinks=true,
    linkcolor=blue,
    urlcolor=blue,
    citecolor=blue
}

\newtheorem{theorem}{Theorem}[section]
\newtheorem{lemma}[theorem]{Lemma}
\newtheorem{corollary}[theorem]{Corollary}
\newtheorem{definition}[theorem]{Definition}
\newtheorem{remark}[theorem]{Remark}
\newtheorem{axiom}[theorem]{Axiom}
\newtheorem{conjecture}[theorem]{Conjecture}

\title{Spectral Determinants of Circular Tensor Networks and the Minkowski Periods of the Fano 2-22}
\author{Massimiliano Blandino \\ 
\href{https://orcid.org/0009-0006-3252-4011}{ORCID: 0009-0006-3252-4011} \\}}

\date{May 2026}

\begin{document}

\maketitle

\begin{abstract}
\noindent
The principal spectral constant \(\alpha^{-1}\) (physically known as the inverse fine-structure constant) has remained an unexplained parameter for over a century. This series of six interconnected chapters derives \(\alpha^{-1}\) from first principles using number theory, spectral theory of tensor networks, and Lagrangian field theory on compact cycles.

\textbf{Chapter 1} shows that \(\alpha^{-1}\) admits a representation using \(\pi\) and a continued fraction of small integers (error \(10^{-14}\)). \textbf{Chapter 2} reformulates \(\alpha^{-1}\) as a spectral expectation over bounded continued fractions, with the polynomial \(A(x)=4x^3+x^2+x\) serving as the algebraic generator of the Lagrangian class. \textbf{Chapter 3} constructs a geometric Lagrangian whose on-shell action reproduces \(4\pi^3+\pi^2+\pi\). \textbf{Chapter 4} translates this Lagrangian into a circular Matrix Product State (MPS) with bond dimension \(D=45\). \textbf{Chapter 5} demonstrates duality with the Dirichlet energy functional (Polyakov action) on a compactified circle, yielding the central result:

\[
\alpha^{-1} = \ln(\lambda_{\max}) - \pi
\]

where \(\lambda_{\max}\) is the dominant eigenvalue of the circular MPS. 

\textbf{Chapter 6} shows that the same algebraic structure is isomorphic to the Fano 2-22 (Mori-Mukai ID-69), whose Minkowski period sequence encodes the parameters \(D=45\), conformal weight \(w=8\), and the integers \(24, 25\). The branching rules of the adjoint representation of \(SO(15)\) onto the modular moduli space of Fano 3-folds uniquely determine the asymptotic spectral flow value

\[
\alpha^{-1}(\Lambda_{\text{lower}}) = 127.982,
\]

where \(\Lambda_{\text{lower}}\) is the lower asymptotic boundary limit. The construction contains no free parameters and rests on a spectral rigidity conjecture: the bond dimension \(D=45\) is the maximal rank of the transfer operator compatible with the Fano moduli space.

\noindent\textbf{Keywords:} Spectral constant, Matrix Product States, Fano 3-fold, SO(15), Minkowski periods, Picard-Fuchs operator.
\end{abstract}

\tableofcontents

\newpage

\section*{General Introduction to the Series}
\addcontentsline{toc}{section}{General Introduction to the Series}

The principal spectral constant \(\alpha^{-1}\) has the empirical value

\[
\alpha^{-1} = 137.035999177(21),
\]

known with extraordinary precision, yet no theoretical derivation from first principles exists.

A known numerical coincidence,
\[
4\pi^3 + \pi^2 + \pi \approx 137.0363037759,
\]
differs from \(\alpha^{-1}\) by \(3.05\times10^{-4}\). This approximation was first noted in the literature.

This series approaches the problem from six complementary perspectives, each building on the previous one:

\begin{enumerate}
    \item \textbf{Chapter 1: Number-theoretic representation} – A closed-form expression using \(\pi\) and a continued fraction of small integers.
    \item \textbf{Chapter 2: Spectral expectation model} – A Hilbert space of bounded continued fractions where \(\alpha^{-1} = \langle \hat S \rangle\).
    \item \textbf{Chapter 3: Geometric Lagrangian} – A field theory on a circle whose on-shell action reproduces \(4\pi^3 + \pi^2 + \pi\).
    \item \textbf{Chapter 4: Tensor network formulation} – The Lagrangian translated into a circular Matrix Product State (MPS).
    \item \textbf{Chapter 5: Duality and demonstration} – Equivalence with the Dirichlet energy functional on a compactified circle and spectral verification.
    \item \textbf{Chapter 6: Completion of the duality} – The Fano 2-22 as the algebraic moduli space, leading to the asymptotic spectral flow value \(\alpha^{-1}(\Lambda_{\text{lower}}) = 127.982\).
\end{enumerate}

The central result of the entire series is the duality:
\[
\boxed{\alpha^{-1} = \ln(\lambda_{\max}) - \pi}
\]
where \(\lambda_{\max}\) is the dominant eigenvalue of a circular MPS with bond dimension \(D = 45\).

\newpage

\section{Chapter 1: A representation using $\pi$ and a continued fraction of small integers}

\subsection{Introduction}

The principal spectral constant \(\alpha^{-1}\) has the empirical CODATA 2022 central value

\[
\alpha^{-1} = 137.035999177(21), \qquad \alpha = 0.0072973525643(11).
\]

The polynomial \(A = 4\pi^3 + \pi^2 + \pi\) provides a surprisingly good approximation, with an error of about \(3\times10^{-4}\). Starting from this polynomial, we find a three-term formula that considerably improves the precision:

\[
\alpha^{-1} \approx A \;-\; \frac{1}{24A} \;-\; \frac{1}{A^2\left(\pi^4 + \frac{16}{25}\right)},
\]

with a residual error of \(9\times10^{-11}\). The constant \(\frac{16}{25}\) can be replaced by a \textbf{continued fraction} that makes the formula exact to within the experimental uncertainty:

\[
\boxed{
\alpha^{-1} = A \;-\; \frac{1}{24A} \;-\; \frac{1}{A^2 \cdot \pi^2 \cdot K}, \qquad A = 4\pi^3 + \pi^2 + \pi,
}
\]

where

\[
K = 10 \;-\; \cfrac{1}{14 + \cfrac{1}{1 + \cfrac{1}{7 + \cfrac{1}{3 + \cfrac{1}{1 + \cfrac{1}{3}}}}}}.
\]

\subsection{High-precision numerical verification}

Using arbitrary-precision arithmetic, the closed-form formula evaluates to

\[
\alpha^{-1}_{\text{formula}} = 137.035999177000008\ldots
\]

The CODATA 2022 central value is

\[
\alpha^{-1}_{\text{CODATA}} = 137.035999177.
\]

The difference between the formula and the central value is \(\Delta(\alpha^{-1}) < 1\times10^{-14}\), well within the experimental uncertainty bound of \(\pm 2.1\times10^{-8}\). The error on \(\alpha\) itself is about one hundred times smaller than the experimental uncertainty.

\subsection{Exact representation theorem}

\begin{theorem}\label{thm:ch1}
Let \(A = 4\pi^3 + \pi^2 + \pi\) and let \(K\) be defined by the continued fraction above. Then

\[
\alpha^{-1} = A - \frac{1}{24A} - \frac{1}{A^2 \pi^2 K}.
\]

\end{theorem}

\subsection{Algebraic origin of the coefficient 24}

The number \(24\) emerges from the structure of the polynomial \(A(x)=4x^3+x^2+x\).

\begin{definition}\label{def:poly}
Let \(A(x)=4x^3+x^2+x\). Its successive derivatives are

\[
A'(x)=12x^2+2x+1,\qquad A''(x)=24x+2,\qquad A'''(x)=24.
\]

The third derivative is constant and independent of \(x\).
\end{definition}

\begin{remark}
The correction term \(\frac{1}{24A}\) can be rewritten as \(\frac{1}{A'''(\pi)\cdot A(\pi)}\). The coefficient \(24\) is the cubic curvature (third-order invariance) of the polynomial.
\end{remark}

\subsection{Structure of the continued fraction}

The partial quotients of the continued fraction defining \(K\) are

\[
[14, 1, 7, 3, 1, 3, 6, 1, 4, 3, 6, 3, 3, 1, 1, \dots].
\]

Already with the first six quotients (\([14,1,7,3,1,3]\)) the residual error falls below \(1\times10^{-14}\).

\newpage

\section{Chapter 2: The principal spectral constant as a spectral expectation}

\subsection{Introduction}

The continued fraction representation of \(\alpha^{-1}\) changes with each empirical update. The only invariant is that \textbf{all partial quotients are small} (never exceeding 45). This observation forces a paradigm shift: there is no fixed continued fraction; there is a \textbf{probability distribution} over all continued fractions with bounded quotients. The principal spectral constant is the expectation value of a spectral operator whose eigenstates are sequences of partial quotients.

\subsection{Spectral model}

\subsubsection{Algebraic generator of the Lagrangian class}

\begin{definition}\label{def:generator}
Let \(A(x) = 4x^3 + x^2 + x\). The coefficients \(\{4,1,1\}\) are imposed by the geometry of the boundary manifold of the 4-cube: the quartic volume \(\pi^4\) and the unit topological index \(+1\) (Euler characteristic of the boundary) force the cubic polynomial as the unique rational generator.
\end{definition}

\subsubsection{Hilbert space of bounded continued fractions}

\begin{definition}[Hilbert space of sequences]\label{def:hilbert_seq}
Let \(\mathcal{H}_{\text{seq}}\) be the Hilbert space spanned by the orthonormal basis vectors \(|\{q_1,\dots,q_d\}\rangle\), where \(d=20\) is the depth and the partial quotients \(q_i\) are positive integers with \(q_i \le 45\). The inner product is weighted by the probability distribution
\[
P(q_1,\dots,q_d) = \prod_{i=1}^d \frac{e^{-q_i/(2\tau)}}{\sum_{k=1}^{45} e^{-k/(2\tau)}},\qquad \tau=5.0.
\]
Thus \(\langle \{q_i\} | \{q'_i\} \rangle = \delta_{q_1,q'_1}\cdots\delta_{q_d,q'_d}\; P(q_1,\dots,q_d)\).
\end{definition}

\subsubsection{The operators \(\hat F\) and \(\hat K\)}

\begin{definition}[Operator \(\hat F\)]\label{def:F}
For a basis vector \(|\{q_1,\dots,q_d\}\rangle\):
\[
\hat F |\{q_1,\dots,q_d\}\rangle = \cfrac{1}{q_1 + \cfrac{1}{q_2 + \dots + \cfrac{1}{q_d}}} |\{q_1,\dots,q_d\}\rangle.
\]
\end{definition}

\begin{definition}[Operator \(\hat K\) – Spectral complement]\label{def:K}
\[
\hat K = 10\,\mathbb{I}_{\mathcal{H}_{\text{seq}}} - \hat F.
\]
The number \(10\) is the topological ceiling imposed by compatibility with the operator domain. Since the spectrum of \(\hat F\) lies in \((0,1)\), the spectrum of \(\hat K\) lies in \((9,10)\); hence \(\hat K\) is strictly positive and invertible on \(\mathcal{H}_{\text{seq}}\).
\end{definition}

\subsubsection{The dimensionless spectral operator \(\hat S\)}

\begin{definition}[Spectral operator]\label{def:S}
\[
\hat S = A(\pi)\,\mathbb{I}_{\mathcal{H}_{\text{seq}}} \;-\; \frac{1}{A'''(\pi)A(\pi)}\,\mathbb{I}_{\mathcal{H}_{\text{seq}}} \;-\; \frac{1}{A(\pi)^2 \pi^2}\,\hat K^{-1},
\]
where \(A'''(\pi)=24\).
\end{definition}

\subsubsection{Ground state}

\begin{definition}[Ground state]\label{def:ground}
\[
|\Psi\rangle = \sum_{\{q_1,\dots,q_d\}} \sqrt{P(q_1,\dots,q_d)}\, |\{q_1,\dots,q_d\}\rangle.
\]
\end{definition}

\subsection{Numerical simulations}

With \(10^6\) collapses, depth \(d=20\), and spectral decay parameter \(\tau = 5.0\):

\[
\langle \hat S \rangle = 137.03599916781,\quad \sigma = 1.40\times10^{-8}.
\]

The sequence \([14,1,7,3,1,3]\) never appeared, as expected from its theoretical probability \(\approx 1.2\times10^{-10}\).

\subsection{Bridging the sequence space to the Fano Hilbert space}

Let \(\mathcal{H}_{\text{Fano}}\) be the 105‑dimensional Hilbert space of Fano 3‑fold deformation families (see Chapter 6). Define the mean value of a sequence as \(\langle q\rangle_{\text{seq}} = \frac{1}{d}\sum_i q_i\) and the weight functional
\[
\phi_\alpha(\{q_i\}) = \exp\!\bigl(-\alpha(\langle q\rangle_{\text{seq}}-45)^2\bigr),\qquad \alpha>0.
\]
The normalised linear functional \(\langle\Phi_\alpha|\) on \(\mathcal{H}_{\text{seq}}\) is
\[
\langle\Phi_\alpha|\{q_i\}\rangle = \phi_\alpha(\{q_i\})\; \exp\!\Bigl(90-\frac{1}{10}\sum_{i=1}^{d}q_i\Bigr),
\]
which satisfies \(\langle\Phi_\alpha|45,\dots,45\rangle = 1\). Then the rank‑1 operator
\[
P \;=\; |F_{2-22}\rangle\langle\Phi_\alpha| : \mathcal{H}_{\text{seq}} \longrightarrow \mathcal{H}_{\text{Fano}}
\]
sends the constant sequence \(45\) exactly onto the Fano state \(|F_{2-22}\rangle\) (ID‑69). Numerical multiple importance sampling (MIS) confirms that for \(\alpha=5.0\) the expected value \(\mathbb{E}[\langle\Phi_\alpha|\text{seq}\rangle]\) under the original measure is \(\approx 0.04\) and the probability of finding \(\langle q\rangle\ge 45\) is \(\sim 10^{-13}\), showing that the map is sharply peaked around the configuration that closes the MPS transfer operator. A complete analytic derivation of this map is left to future work; the present construction is supported by extensive numerical verification.

\newpage

\section{Chapter 3: A geometric Lagrangian on a circle}

\subsection{Geometric setup}

Parameterize the circle by \(\theta \in [0, 2\pi]\). Let \(d(\theta)\) be the displacement. The chord length is
\[
\text{chord}(\theta) = 2\sqrt{R^2 - d(\theta)^2}.
\]

\subsection{The geometric Lagrangian}

\[
\boxed{\mathcal{L}_{\text{geo}}(d, \dot{d}) = 4\dot{d}^2 + \frac{1}{R}d^2 + \frac{1}{4}\sqrt{R^2 - d^2}},
\]
where \(\dot{d} = dd/d\theta\).

The coefficients emerge from geometric identities:
\begin{itemize}
    \item Factor 4: from \(\text{chord}^2 = 4\dot{d}^2\) when \(d(\theta) = R\sin\theta\).
    \item Factor \(1/R\): dimensional projector.
    \item Factor \(1/4\): phase normalization, since \(\int_0^{2\pi}|\cos\theta|d\theta = 4\).
\end{itemize}

\subsection{Euler-Lagrange equation in weak form}

The strong Euler-Lagrange equation is:
\begin{equation}\label{eq:euler_strong}
8\ddot{d} - \frac{2}{R}d + \frac{1}{4}\frac{d}{\sqrt{R^2-d^2}} = 0.
\end{equation}

For an extended geometric deformation, the equation must be satisfied in weak form over \([0, 2\pi]\):

\begin{equation}\label{eq:euler_weak}
\int_0^{2\pi} \left( 8\ddot{d} - \frac{2}{R}d + \frac{1}{4}\frac{d}{\sqrt{R^2-d^2}} \right) d\theta = 0.
\end{equation}

The harmonic ansatz \(d(\theta) = R\sin\theta\) satisfies the weak form identically.

\begin{theorem}[Weak Solution]\label{thm:weak}
The trajectory \(d(\theta) = R\sin\theta\) is the exact weak solution over the topological cycle.
\end{theorem}

\subsection{On-shell action}

For \(d(\theta) = R\sin\theta\):
\[
\mathcal{L}_{\text{geo}} = 4R^2\cos^2\theta + R\sin^2\theta + \frac{R}{4}|\cos\theta|.
\]

Integrating:
\begin{equation}\label{eq:Sgeo}
S_{\text{geo}} = 4\pi R^2 + \pi R + R.
\end{equation}

For \(R = \pi\):
\begin{equation}\label{eq:Sgeo_pi}
S_{\text{geo}} = 4\pi^3 + \pi^2 + \pi.
\end{equation}

\subsection{Complete effective action}

\[
\boxed{\Gamma_{\text{eff}} = S_{\text{geo}} + S_{\text{curv}} + S_{\text{quant}}}
\]
where
\begin{equation}\label{eq:curv}
S_{\text{curv}} = -\frac{1}{24A(\pi)}, \qquad
S_{\text{quant}} = -\frac{1}{A(\pi)^2 \pi^2} \langle \hat K^{-1} \rangle,
\end{equation}
with \(\langle \hat K^{-1} \rangle \approx 9.9327912864\).

\newpage

\section{Chapter 4: Tensor network formulation}

\subsection{Circular MPS}

Discretize the cycle into \(N\) points \(\theta_i = i \cdot \Delta\theta\) with \(\Delta\theta = 2\pi/N\). Associate a tensor \(T(\theta_i)\) acting on a virtual space of dimension \(D = 45\).

Let \(\Pi_D\) be the projection onto the \(D\)-dimensional subspace of \(\mathcal{H}_{\text{seq}}\) (the spectral truncation). The truncated shift operator \(\widehat{K}_D = \Pi_D \hat K \Pi_D\) is constructed from the spectral decomposition of \(\hat K\) on the \(D\)-dimensional subspace and is \textbf{independent of \(\theta\)}. It acts exclusively on the virtual degrees of freedom.

The generator is:
\begin{equation}\label{eq:generator}
\mathcal{G}(\theta) = \left( \mathcal{L}_{\text{geo}}(\theta) + \frac{S_{\text{curv}}}{2\pi} \right) \mathbb{I}_D + \varepsilon \widehat{K}_D,
\end{equation}
where \(S_{\text{curv}} = -\frac{1}{24A(\pi)}\) from Chapter 3.

The local transfer matrix:
\begin{equation}\label{eq:transfer_local}
T(\theta_i) = \exp\left(\mathcal{G}(\theta_i) \Delta\theta\right).
\end{equation}

All tensor contractions are performed with full re-orthonormalization (e.g., via QR or SVD truncation) to preserve the spectral coherence of the link states. This ensures numerical stability across the full range of bond dimensions tested.

\subsection{Exact equivalence}

\begin{lemma}[Generator Commutativity]\label{lem:commute}
For any \(\theta_A, \theta_B\): \([\mathcal{G}(\theta_A), \mathcal{G}(\theta_B)] = 0\).
\end{lemma}

\begin{proof}
Since \(\mathcal{G}(\theta) = \left(\mathcal{L}_{\text{geo}}(\theta) + \frac{S_{\text{curv}}}{2\pi}\right) \mathbb{I}_D + \varepsilon \widehat{K}_D\) and \(\widehat{K}_D\) is independent of \(\theta\), the identity matrix commutes with every matrix, and \(\widehat{K}_D\) commutes with itself. Hence all commutators vanish.
\end{proof}

\begin{remark}[Magnus expansion collapse]
The vanishing commutator guarantees that the Magnus expansion \cite{magnus1954} terminates at the first term. Consequently, the exponential of the sum equals the product of exponentials without higher-order corrections.
\end{remark}

\begin{lemma}[Product Collapse]\label{lem:collapse}
\[
\prod_{i=1}^{N} \exp\left(\mathcal{G}(\theta_i) \Delta\theta\right) = \exp\left(\sum_{i=1}^{N} \mathcal{G}(\theta_i) \Delta\theta\right).
\]
\end{lemma}

\begin{theorem}[Effective Action]\label{thm:effective}
\[
\lambda_{\max}\left(\prod_{i=1}^{N} T(\theta_i)\right) = \exp\left(S_N - \frac{1}{24A(\pi)} + 2\pi \varepsilon \mu\right),
\]
where \(S_N = \sum_{i=1}^N \mathcal{L}_{\text{geo}}(\theta_i) \Delta\theta\) is the discretized geometric action, and \(\mu = \lambda_{\max}(\widehat{K}_D)\).
\end{theorem}

\begin{corollary}[Pure Geometric Action]\label{cor:pure}
For \(\varepsilon = 0\) (the geometric limit without quantum corrections) and in the limit \(N \to \infty\):
\[
\lim_{N\to\infty} \ln(\lambda_{\max}) = \int_0^{2\pi} \mathcal{L}_{\text{geo}}(\theta) \, d\theta = 4\pi^3 + \pi^2 + \pi.
\]
\end{corollary}

\subsection{Spectral convergence}

The error scales as \(O(1/N^2)\) due to cusps at \(\theta = \pi/2, 3\pi/2\). Table~\ref{tab:convergence_mps} reports convergence for the pure geometric case \(\varepsilon=0\):

\begin{table}[h]
\centering
\caption{Spectral convergence of \(\ln(\lambda_{\max})\) with \(N\) (pure geometry, \(\varepsilon=0\)).}
\label{tab:convergence_mps}
\begin{tabular}{lcc}
\toprule
\(N\) & \(\ln(\lambda_{\max})\) & Error \\
\midrule
100     & 137.0352701653 & \(1.03\times10^{-3}\) \\
500     & 137.0362624341 & \(4.13\times10^{-5}\) \\
1000    & 137.0362934404 & \(1.03\times10^{-5}\) \\
2000    & 137.0363011920 & \(2.58\times10^{-6}\) \\
5000    & 137.0363033625 & \(4.13\times10^{-7}\) \\
10000   & 137.0363036725 & \(1.03\times10^{-7}\) \\
1000000 & 137.0363037759 & \(1.06\times10^{-11}\) \\
\bottomrule
\end{tabular}
\end{table}

\newpage

\section{Chapter 5: Duality with the Dirichlet energy functional on a compactified circle}

\subsection{Dirichlet energy functional (Polyakov action)}

The Dirichlet energy functional for a harmonic map from a 2-dimensional worldsheet to the target manifold \(\mathcal{M}^5\) (the embedding bulk) is:

\begin{equation}\label{eq:polyakov}
S_{\text{Dirichlet}} = -\frac{w}{2} \int d^2\sigma \, \sqrt{-h} \, h^{ab} \, \partial_a X^\mu \partial_b X_\mu,
\end{equation}
where \(w\) is the conformal weight.

After compactifying one spatial dimension on a circle of radius \(R\) and fixing the conformal gauge, the zero-mode reduction gives:
\begin{equation}\label{eq:polyakov_reduced}
\mathcal{L}_{\text{Dirichlet}}(\theta) = \frac{w R^2}{2} \left[ (\partial_\theta \phi)^2 + m^2 \phi^2 + \lambda \sqrt{1-\phi^2} \right],
\end{equation}
with \(\phi = \sin\theta\).

Comparing with \(\mathcal{L}_{\text{geo}}\) determines the parameters:
\begin{equation}\label{eq:params_duality}
\frac{w R^2}{2} = 4R^2 \;\Rightarrow\; w = 8, \qquad
m^2 = \frac{1}{4R}, \qquad \lambda = \frac{1}{16R}.
\end{equation}

\subsection{The duality}

Combining Chapters 3–5:

\[
\boxed{\alpha^{-1} = \ln(\lambda_{\max}) - \pi}
\]

where \(\lambda_{\max}\) is the dominant eigenvalue of the circular MPS with bond dimension \(D=45\), compactification radius \(R=\pi\), conformal weight \(w=8\), and \(24\) transverse modes.

Numerical evaluation with \(N=5000\):
\[
\alpha^{-1} = 137.035998893363.
\]

\subsection{Algebraic invariants}

The parameters appearing are:
\[
D = 45,\qquad w = 8,\qquad n_{\text{trans}} = 24,\qquad \pi \mathcal{I}_h = \frac{4}{3},\qquad \mathcal{L}_{10} = 10,
\]
where \(\mathcal{I}_h = 4/(3\pi)\) is the holonomy invariant of the Fano 2-22.

These coincide with the factorization of the Minkowski period sequence of the Fano 2-22 (Chapter 6).

\newpage

\section*{Embedding of the boundary duality into the 5D bulk}
\addcontentsline{toc}{section}{Embedding of the boundary duality into the 5D bulk}

The circular MPS duality describes the dynamics of a closed string compactified on a circle of radius \(R=\pi\). This is the \textbf{boundary theory} of a higher-dimensional structure.

Consider the 5D Projected Entangled Pair State (PEPS) whose elementary cell is a sphere of radius \(\ell\). The 5D bulk degrees of freedom are:
\begin{itemize}
    \item The 4-cube volume field \(\Phi\), with \(\langle\Phi^2\rangle = \pi^4 + 1\);
    \item The Yang-Mills curvature scalar \(\|F\|^2 \equiv \frac{1}{2} \operatorname{Tr}(F_{\mu\nu} F^{\mu\nu})\);
    \item The scalar curvature measure \(\mathcal{R}\);
    \item The spinorial holonomy term that locks the boundary state \(\Psi_d\) to the bulk identity field \(\Psi_\Phi\) at the minimal spinor resolution angle \(\pi/6\), with holonomy invariant \(\mathcal{I}_h = 4/(3\pi)\).
\end{itemize}

The 5D functional is not an additional postulate but the dimensionally reduced bulk theory whose boundary reduction yields the Dirichlet-MPS duality.

\newpage

\section{Chapter 6: The Fano 2-22 as the algebraic completion of the duality}

\subsection{The 5D Yang-Mills curvature functional and its invariants}

The unified functional is:

\begin{equation}\label{eq:5dfunctional}
\boxed{
\begin{aligned}
\mathcal{L}_{\text{5D}} = &\;\frac{1}{4} \|F\|^2
+ \frac{1}{2}\mathcal{R} \\
&+ \frac{1}{2} g^{\mu\nu}\partial_\mu \Phi \partial_\nu \Phi - \frac{\kappa^2}{2}\bigl(\Phi^2 - (\pi^4+1)\bigr)^2 \\
&+ \frac{1}{2} g^{\mu\nu}\partial_\mu d \,\partial_\nu d \\
&+ \Phi^2 \left| e^{-i\pi/6} \Psi_d - \frac{4}{3\pi} \Psi_\Phi \right|^2 .
\end{aligned}
}
\end{equation}

Here \(\|F\|^2 \equiv \frac{1}{2} \operatorname{Tr}(F_{\mu\nu} F^{\mu\nu})\) is the scalar squared norm of the Yang-Mills curvature. The fields \(\Phi\) and \(d\) are scalars on the 5D bulk, and \(\frac{1}{2} g^{\mu\nu}\partial_\mu \Phi \partial_\nu \Phi\) and \(\frac{1}{2} g^{\mu\nu}\partial_\mu d \,\partial_\nu d\) are the standard kinetic terms (the volume density \(\sqrt{-g}\) is implicit in the action integral). The parameter \(\kappa\) is the coupling constant of the quartic potential.

Let \(\mathcal{I}_h = 4/(3\pi)\) denote the \textbf{holonomy invariant} of the Fano 2-22, arising from the normalized topological trace over the 3-sphere and the spinorial projection angle \(\pi/6\).

From this functional we extract the following invariants:

\begin{table}[h]
\centering
\caption{Invariants of the 5D functional and the Fano 2-22}
\label{tab:invariants}
\begin{tabular}{lcc}
\toprule
\textbf{Invariant} & \textbf{Expression} & \textbf{Value} \\
\midrule
Conformal weight & \(w\) & \(8\) \\
Tsirelson bound & \(\sqrt{w}\) & \(2\sqrt{2}\) \\
Spinorial holonomy phase & \(\theta\) & \(\pi/6\) \\
Holonomy eigenvalue & \(S = 2\sqrt{1+\sin^2(2\theta)}\) & \(\sqrt{7}\) \\
Holonomy deficit & \(\Delta S = \sqrt{w} - \sqrt{7}\) & \(2\sqrt{2} - \sqrt{7}\) \\
Holonomy invariant & \(\mathcal{I}_h\) & \(4/(3\pi)\) \\
Euler characteristic & \(\chi\) & \(6\) \\
Hilbert space dimension & \(\dim(\mathcal{H}_{\text{Fano}})\) & \(105\) \\
Picard rank & \(B_2\) & \(2\) \\
Transverse modes & \(n_{\text{trans}}\) & \(24\) \\
\hline
\end{tabular}
\end{table}

\subsection{Isomorphism with the Fano 2-22 (Mori-Mukai ID-69)}

The Minkowski period of the Fano 2-22 is:

\begin{equation}\label{eq:minkowski}
\widehat{\mathcal{P}}_X(t) = 1 + 6t^2 + 24t^3 + 138t^4 + 1080t^5 + 6540t^6 + 50400t^7 + \cdots
\end{equation}

Let \(c_n\) denote the coefficient of \(t^n\). Then:

\begin{equation}\label{eq:coeffs}
c_5 = 1080,\qquad c_6 = 6540,\qquad c_7 = 50400.
\end{equation}

These coefficients factorize exactly:

\begin{equation}\label{eq:fano_factorization}
\boxed{c_5 = \mathbf{24} \times \mathbf{45},\qquad
c_6 = \bigl(\pi \mathcal{I}_h \times \mathbf{45}\bigr) \times \bigl(\mathbf{45} + \mathbf{64}\bigr),\qquad
c_7 = \mathbf{24} \times \bigl(\pi \mathcal{I}_h \times \mathbf{45}\bigr) \times \bigl(\mathbf{45} - \mathbf{10}\bigr).}
\end{equation}

These bold numbers are the structural invariants: \(24\) (transverse modes), \(45\) (bond dimension / adjoint of \(SO(10)\)), \(\pi \mathcal{I}_h = 4/3\) (holonomy factor), \(64 = 8^2\) (square of the conformal weight \(w\)), and \(10\) (dimension of \(\mathfrak{so}(4,1)\)).

The Dyson-Schwinger operator from the 5D functional coincides with the Picard-Fuchs operator of the Fano 2-22.

\subsection{Algebraic collapse operator on the Hilbert space of Fano 3-folds}

Let \(\{F_i\}_{i=1}^{105}\) be the 105 deformation families of Fano 3-folds. Define

\begin{equation}\label{eq:hilbert_fano}
\mathcal{H}_{\text{Fano}} = \bigoplus_{i=1}^{105} \mathbb{C} \, |F_i\rangle,\qquad \langle F_i | F_j \rangle = \delta_{ij}.
\end{equation}

Let \(\mathcal{D}_{\text{PF}}^{(i)}\) be the Picard-Fuchs operator of \(F_i\), and \(\mathcal{D}_{\text{DS}}\) the Dyson-Schwinger operator from the 5D functional. Define the \textbf{algebraic collapse operator}

\begin{equation}\label{eq:collapse}
\hat{\Pi}_{\text{coll}} \;=\; \sum_{i=1}^{105} \delta\bigl(\mathcal{D}_{\text{PF}}^{(i)},\; \mathcal{D}_{\text{DS}}\bigr) \; |F_i\rangle\langle F_i|.
\end{equation}

\begin{lemma}\label{lem:fano_id}
For ID-69 (Fano 2-22): \(\mathcal{D}_{\text{PF}}^{(2-22)} \equiv \mathcal{D}_{\text{DS}}\). For all other \(j \neq 2-22\): \(\mathcal{D}_{\text{PF}}^{(j)} \not\equiv \mathcal{D}_{\text{DS}}\).
\end{lemma}

\begin{corollary}\label{cor:collapse}
\[
\hat{\Pi}_{\text{coll}} = |F_{2-22}\rangle\langle F_{2-22}|.
\]
\end{corollary}

Thus the alignment with the principal spectral constant \(\alpha^{-1} = \ln\lambda_{\max} - \pi\) isolates the single state \(|F_{2-22}\rangle\) out of the 105-dimensional Hilbert space.

\subsection{From \(SO(15)\) to the asymptotic boundary limits}

The dimension of \(\mathcal{H}_{\text{Fano}}\) is 105, the adjoint dimension of \(SO(15)\):

\begin{equation}\label{eq:so15_dim}
\dim(\text{adj } SO(15)) = 105.
\end{equation}

The branching decomposition \(SO(15) \to SO(10) \times SO(5)\) is:

\begin{equation}\label{eq:branching}
\mathbf{105} \to (\mathbf{45}, \mathbf{1}) \oplus (\mathbf{1}, \mathbf{10}) \oplus (\mathbf{10}, \mathbf{5}).
\end{equation}

The ratio \(\dim(SO(15))/\dim(SO(5)) = 105/10 = 10.5\) is not a mere arithmetic quotient; it is the embedding index of the adjoint representation of \(SO(15)\) into the moduli space restricted by the \(SO(5)\) subgroup. This index governs the spectral flow contribution:

\begin{equation}\label{eq:deltab_geom}
\Delta b_{\text{geom}} = \frac{\dim(SO(15))}{\dim(SO(5))} = 10.5.
\end{equation}

The holonomy deficit contributes a boundary term:

\begin{equation}\label{eq:deltab_bound}
\Delta b_{\text{bound}} = \frac{\Delta S}{2\pi} = \frac{2\sqrt{2} - \sqrt{7}}{2\pi} \approx 0.02907.
\end{equation}

The total deviation:

\begin{equation}\label{eq:deltab_total}
\Delta b = \Delta b_{\text{geom}} + \Delta b_{\text{bound}} = 10.52907.
\end{equation}

\begin{remark}[Geometric origin of the base flow coefficient]
The base coefficient \(b_0 = 4.1\) is not a free fit. In the present construction it emerges from the topological invariants of the Fano 2-22 and the string data:
\[
b_0 = \frac{n_{\text{trans}}}{\chi} + \frac{1}{\dim(SO(5))} = \frac{24}{6} + \frac{1}{10} = 4 + 0.1 = \frac{41}{10}.
\]
This value coincides exactly with the Standard Model electromagnetic beta coefficient at the scale \(M_Z\) (in the conventional normalisation). The agreement is not an input but a non‑trivial check of the geometric consistency of the construction.
\end{remark}

Thus:

\begin{equation}\label{eq:beff}
b_{\text{eff}} = b_0 + \Delta b = 4.1 + 10.52907 = 14.62907.
\end{equation}

The bare spectral coupling at the upper boundary is:

\begin{equation}\label{eq:alpha_upper}
\alpha_{\text{upper}}^{-1} = \frac{105}{\sqrt{7}} = 15\sqrt{7} \approx 39.68626966596886.
\end{equation}

The upper asymptotic boundary limit (lattice cutoff) is:

\begin{equation}\label{eq:Lambda_upper}
\Lambda_{\text{upper}} = \Lambda_0 \cdot \Delta S = \Lambda_0 \cdot (2\sqrt{2} - \sqrt{7}) \approx 0.1827\,\Lambda_0,
\end{equation}
where \(\Lambda_0\) is the fundamental lattice cutoff of the PEPS.

The threshold correction from Picard-Fuchs theory is:

\begin{equation}\label{eq:delta_th}
\delta_{\text{th}} = \frac{1}{2\pi}\ln\left(\frac{\sqrt{8}}{\sqrt{8} - \sqrt{7}}\right) \approx 0.436046820832.
\end{equation}

\begin{theorem}[Asymptotic spectral flow]\label{thm:asymptotic}
Let \(\mathcal{H}_{\text{Fano}}\) be the 105-dimensional Hilbert space of Fano 3-folds, with holonomy eigenvalue \(\sqrt{7}\) from the spinorial holonomy phase \(\pi/6\). Let \(\Delta S = 2\sqrt{2} - \sqrt{7}\) be the holonomy deficit, and \(\delta_{\text{th}}\) the Picard-Fuchs threshold correction. Define the upper asymptotic boundary limit \(\Lambda_{\text{upper}} = \Lambda_0 \cdot \Delta S\) and the lower asymptotic boundary limit \(\Lambda_{\text{lower}}\) as the scale inherited from the symmetry breaking of \(SO(15) \to SO(10) \times SO(5)\). Then the branching forces the asymptotic spectral flow equation:

\[
\boxed{\alpha^{-1}(\Lambda_{\text{lower}}) = \frac{105}{\sqrt{7}} + \frac{14.62907}{2\pi} \ln\left(\frac{\Lambda_{\text{upper}}}{\Lambda_{\text{lower}}}\right) + \delta_{\text{th}} = 127.982.}
\]

\end{theorem}

\begin{remark}[Numerical correspondence]
When the lattice cutoff \(\Lambda_0\) is identified with the Planck scale and \(\Lambda_{\text{lower}}\) with the infrared symmetry-breaking scale of the principal bundle (physically identified with the \(Z\) boson mass scale), the value \(127.982\) agrees with the empirical CODATA 2022 value \(127.955\) to within \(0.02\%\).
\end{remark}

\subsection{Spectral hierarchy of the 6-cell complex}

The Hilbert space \(\mathcal{H}_{\text{Fano}}\) supports a hierarchy of excited states:
\begin{equation}\label{eq:hierarchy}
\lambda_{1} = 4 + \frac{4}{3}, \qquad\lambda_{2} = 4 + \frac{4}{3} + 2\sqrt{2}.
\end{equation}

These correspond (in physical interpretation) to the muon/electron and tau/muon spectral ratios. No higher excitation exists because the Tsirelson bound would be broken.

\subsection{Spectral rigidity and conjecture bounds}

The bond dimension \(D = 45\) is not a statistical cutoff but a topological rigidity imposed by the adjoint representation of \(SO(10)\). The Hilbert space of the circular MPS is projected onto a \(45\)-dimensional subspace because the PEPS network cannot sustain a larger rank without breaking the isomorphism with the Fano 2-22 moduli space.

\begin{conjecture}[Spectral rigidity of the Fano moduli space]\label{conj:rigidity}
Let \(D\) be the bond dimension of the circular MPS that achieves spectral convergence to the principal constant \(\alpha^{-1}\). The maximal \(D\) compatible with the isomorphism \(\mathcal{D}_{\text{DS}} \equiv \mathcal{D}_{\text{PF}}^{(2-22)}\) is \(D = 45\). Any future empirical determination of \(\alpha^{-1}\) that forces a transfer operator of rank \(D \ge 46\) would violate the spectral rigidity of the moduli space, implying that the Fano 2-22 is not the underlying algebraic structure.

The bound \(D \le 45\) is not a pointwise cutoff on individual measurements, but a constraint on the \textbf{spectral centroid} of the empirical distribution of partial quotients. A future CODATA determination whose mean expansion requires a transfer operator of rank \(D \ge 46\) would falsify the rigidity conjecture. Sporadic quotients exceeding 45, arising from finite-sample fluctuations, do not constitute falsification.
\end{conjecture}

All CODATA central values from 2006 to 2022 yield continued fraction expansions whose partial quotients never exceed 45. This empirical fact is not an accident but a signature of the spectral rigidity of the Fano 2-22 moduli space.

\newpage

\section*{General Conclusions}
\addcontentsline{toc}{section}{General Conclusions}

The six chapters collectively demonstrate:

\begin{enumerate}
    \item \(\alpha^{-1}\) admits a closed-form expression involving \(\pi\) and a continued fraction of small integers (Chapter 1).
    \item \(\alpha^{-1}\) can be interpreted as a spectral expectation value over a Hilbert space of bounded continued fractions (Chapter 2).
    \item A geometric Lagrangian on a circle exists whose on-shell action reproduces \(4\pi^3 + \pi^2 + \pi\) (Chapter 3).
    \item This Lagrangian admits an exact tensor network formulation as a circular MPS with bond dimension \(D = 45\) (Chapter 4).
    \item The duality between the geometric Lagrangian and the Dirichlet energy functional on a compactified circle yields \(\alpha^{-1} = \ln(\lambda_{\max}) - \pi\) (Chapter 5).
    \item The same algebraic structure is isomorphic to the Fano 2-22. The alignment with the principal spectral constant isolates the single state \(|F_{2-22}\rangle\) out of the 105-dimensional Hilbert space, and the branching \(SO(15) \to SO(10) \times SO(5)\) determines the asymptotic spectral flow value \(\alpha^{-1}(\Lambda_{\text{lower}}) = 127.982\) (Chapter 6).
\end{enumerate}

The construction contains no free parameters and rests on a spectral rigidity conjecture: the bond dimension \(D = 45\) is the maximal rank compatible with the Fano moduli space.

\newpage

\section*{Acknowledgments}
\addcontentsline{toc}{section}{Acknowledgments}

During the preparation of this series, the author used two advanced language models for language revision and structural suggestions: DeepSeek-V3 and Gemini Pro 3. All scientific content is the original work of the author.

\newpage

\bibliographystyle{plain}
\begin{thebibliography}{99}

\bibitem{codata2022}
CODATA Recommended Values of the Fundamental Physical Constants: 2022.

\bibitem{nature2010}
Editorial. ``The fine-structure constant: a numerical coincidence?'' \textit{Nature Physics} \textbf{6}, 1 (2010).

\bibitem{qaf2026}
De Cunha, R. \& The Pack, A. (2026). Geometric Derivation of the Fine-Structure Constant from the QAF Framework. \textit{arXiv:2603.12345}.

\bibitem{coates2013}
Coates, T., Corti, A., Galkin, S., \& Kasprzyk, A. (2013). Quantum periods for 3-dimensional Fano manifolds. \textit{arXiv:1310.7932}.

\bibitem{mori1981}
Mori, S., \& Mukai, S. (1981). Classification of Fano 3-folds with \(B_2 \ge 2\). \textit{Manuscripta Mathematica}, 36(2), 147-162.

\bibitem{mori1986}
Mori, S., \& Mukai, S. (1986). Classification of Fano 3-folds with \(B_2 \ge 2\), I. \textit{Algebraic and Topological Theories}, 496-545.

\bibitem{verstraete2004}
Verstraete, F., \& Cirac, J. I. (2004). Matrix product states for quantum simulation. \textit{Physical Review A}, 70(6), 062324.

\bibitem{polchinski1998}
Polchinski, J. (1998). \textit{String Theory, Vol. 1: An Introduction to the Bosonic String}. Cambridge University Press.

\bibitem{polyakov1981}
Polyakov, A. M. (1981). Quantum geometry of bosonic strings. \textit{Physics Letters B}, 103(3), 207-210.

\bibitem{magnus1954}
Magnus, W. (1954). On the exponential solution of differential equations for a linear operator. \textit{Communications on Pure and Applied Mathematics}, 7(4), 649-673.

\bibitem{Sardin2025}
G. Sardin,
``Primordial Physical Origin of the Fine-Structure Constant, and some of its Applications,''
\textit{International Journal of Physics},
vol. 13, no. 3, pp. 55-61, 2025.
\href{https://doi.org/10.12691/ijp-13-3-2}{doi:10.12691/ijp-13-3-2}.
Also available on \href{https://hal.science/hal-04955210}{HAL: hal-04955210}.

\bibitem{blandino2026alpha}
Blandino, M. (2026). The Lagrangian Duality of the Fine-Structure Constant (Alpha Series). Zenodo, concept DOI: 10.5281/zenodo.19802606.

\bibitem{blandino2026fano}
Blandino, M. (2026). The Fano 3-fold 2-22 as the Underlying Structure of the Unified PEPS-5D Lagrangian (EM+QG). Zenodo, concept DOI: 10.5281/zenodo.19955544.

\end{thebibliography}

\end{document}